% =============================================================================
% CAAZ SAR ADC 研究工作总结报告
% 编译: xelatex -> xelatex (两次)
% 忠实转换自 研究报告.md，仅更新数据，保留全部原始框架
% =============================================================================
\documentclass[11pt,a4paper,UTF8]{ctexart}
\usepackage{geometry}\geometry{left=2.4cm,right=2.4cm,top=2.5cm,bottom=2.5cm}
\usepackage{amsmath,amssymb,graphicx,booktabs,array,hyperref,caption,titlesec,fancyhdr}
\usepackage[most]{tcolorbox}
\graphicspath{{figures/}}

\newtcolorbox{prinbox}[1]{enhanced,breakable,
  colback=blue!4!white,colframe=blue!50!black,fonttitle=\bfseries,
  borderline west={3pt}{0pt}{blue!60!black},title={\strut #1}}
\newtcolorbox{findbox}[1]{enhanced,breakable,
  colback=green!3!white,colframe=green!50!black,fonttitle=\bfseries,
  borderline west={3pt}{0pt}{green!60!black},title={\strut #1}}
\newtcolorbox{warnbox}[1]{enhanced,breakable,
  colback=orange!3!white,colframe=orange!60!black,fonttitle=\bfseries,
  borderline west={3pt}{0pt}{orange!80!black},title={\strut #1}}
\newtcolorbox{impbox}[1]{enhanced,breakable,
  colback=red!3!white,colframe=red!60!black,fonttitle=\bfseries,
  borderline west={3pt}{0pt}{red!70!black},title={\strut #1}}

\definecolor{pb}{RGB}{0,82,136}\definecolor{db}{RGB}{0,51,102}
\titleformat{\section}{\Large\bfseries\color{pb}}{\thesection}{1em}{}
\titleformat{\subsection}{\large\bfseries\color{db}}{\thesubsection}{1em}{}
\pagestyle{fancy}\fancyhf{}
\fancyhead[L]{\small\color{pb}CAAZ SAR ADC 研究工作总结报告}
\fancyhead[R]{\small\color{pb}HUST-SAR-2026-004}
\fancyfoot[C]{\thepage}\renewcommand{\headrulewidth}{0.4pt}\setlength{\headheight}{22pt}

\begin{document}
\begin{titlepage}\thispagestyle{empty}
\begin{center}\vspace*{2cm}
{\LARGE\bfseries\color{pb}CAAZ SAR ADC 研究工作总结报告}\\[0.3cm]
{\large --- 电流域 $kT/C$ 噪声消除技术 ---}\\[1cm]
{\small HUST-SAR-2026-004 \hspace{2em} Rev.4.0 \hspace{2em} 2026-05-13}\\[1.5cm]
\begin{tcolorbox}[colback=blue!5!white,colframe=pb,sharp corners,title=\textbf{\large\color{pb}核心发现}]
\small\begin{tabular}{ll}
$\bullet$ & 架构突破：CAAZ 将噪声存储电容置于低阻抗节点，增益与饱和解耦，$A<10\rightarrow A>50$ \\
$\bullet$ & 基线性能：NF=1 全五角 ENOB $>$11.4 bit，TT 角 ENOB=\textbf{11.61 bit}，SNR=\textbf{71.81 dB}，SFDR=\textbf{84.01 dB} \\
$\bullet$ & CAAZ 效率：NF=10 双-$t_d$ 法，TT $\eta$=\textbf{87.9\%}（降噪 8.3$\times$），全角 79.7\%--91.2\% \\
$\bullet$ & 跨批次验证：$\eta$ 偏差 $<$1.4 p.p.，COMPOWER $<$3 $\mu$W，结论稳健 \\
$\bullet$ & 功耗 720 $\mu$W（模拟 65\%，数字 35\%），$C_{AZ}=300$ fF，$t_d=26$ ns 最优 \\
$\bullet$ & 早期 10-MS/s 基线：KT/C 噪声功率压缩 5.53$\times$，SNR 提升 $+$1.21 dB \\
\end{tabular}\end{tcolorbox}
\vspace{1.5cm}
\begin{tabular}{ll}
\textbf{工艺：}TSMC 18RF (1P6M, 1.8 V) & \textbf{仿真器：}Cadence Spectre 23.1 (APS)\\
\textbf{作者：}赵诣, 华中科技大学 集成电路学院 & \textbf{时间跨度：}2026年3月--5月\\
\end{tabular}
\end{center}
\end{titlepage}
\tableofcontents\newpage

% ===== 摘要 =====
\section{摘要}
\label{sec:abstract}

\begin{findbox}{核心发现摘要（2026年5月更新）}
\begin{itemize}
  \item \textbf{架构突破}：CAAZ 电流域架构将噪声存储电容置于低阻抗节点，实现增益与饱和解耦，可用增益从 $A<10$ 提升至 $A>50$
  \item \textbf{10-MS/s 基线性能}：NF=1 条件下，全五角 PVT ENOB 均 $>$11.4 bit，TT 角最佳 ENOB=\textbf{11.61 bit}，SNR=\textbf{71.81 dB}，SFDR=\textbf{84.01 dB}
  \item \textbf{CAAZ 效率量化}：通过 NF=10（100$\times$噪声注入）双 $t_d$ 方法，TT 角消除效率 $\eta$=\textbf{87.9\%}，噪声功率压缩 \textbf{8.3$\times$}
  \item \textbf{PVT 全角覆盖}：$\eta$ 范围 79.7\%--91.2\%，SF 角（85$^\circ$C，slow-N fast-P）为最差设计瓶颈
  \item \textbf{跨批次验证}：双批次 CSV 导出数据 $\eta$ 偏差 $<$1.4 p.p.，COMPOWER 一致性 $<$3 $\mu$W，CAAZ 定性结论稳健
  \item \textbf{功耗 720 $\mu$W}（模拟 65\%，数字 35\%），预放大器 PREWOCM 占 41\% 为最大优化空间
  \item \textbf{失调鲁棒性}：注入失调从 0 mV 增至 5 mV 时，ENOB 仅下降约 0.2 bit
  \item \textbf{早期 10-MS/s 基线}：KT/C 噪声功率压缩 5.53$\times$（81.93\% 能量削减），整机 SNR 从 69.59 dB 提升至 70.80 dB（$+$1.21 dB）
\end{itemize}
\end{findbox}

本报告系统梳理了 2026 年 3 月至 5 月期间，基于电流调整自动调零（CAAZ）技术的 SAR ADC 研究工作。报告从研究背景出发，深入分析 CAAZ 技术核心原理，详细介绍了实验平台搭建过程，总结了关键实验发现。研究包含两个阶段：\textbf{Phase I}（3--4月）完成 10-MS/s 方案验证与噪声分解；\textbf{Phase II}（4--5月）推进至 10-MS/s 全面表征，完成了涵盖 TT/FF/SS/SF/FS 五工艺角的 PVT 全面仿真与 NF=10 CAAZ 效率定量表征，并进行了双批次跨验证。

\section{研究背景与动机}
\label{sec:background}

\begin{findbox}{核心亮点}
SAR ADC 在高精度应用中面临 $kT/C$ 噪声限制。传统增大采样电容方案导致驱动功耗居高不下。噪声消除技术虽可突破此限制，但存在\textbf{增益--饱和矛盾}这一核心瓶颈。电流域架构为突破传统 OOS 技术限制提供了新路径。
\end{findbox}

\subsection{问题陈述与研究前沿分析}

逐次逼近型模数转换器（SAR ADC）在高精度应用中面临严峻的采样噪声限制。在典型设计场景下，采样噪声往往占据整体 ADC 噪声预算的主导地位。传统抑制方法采用增大采样电容的策略——文献 [1] 即采用了高达 15 pF 的采样电容——但这导致前端驱动电路功耗居高不下，严重制约系统能效。

\textbf{噪声消除技术的演进脉络}

2020 年，清华大学孙楠教授团队提出了面向 SAR ADC 的 $kT/C$ 噪声消除技术 [2]。该技术源于 ADI 团队 2014 年的采样保持电路方案 [3]，核心洞察在于：当采样电容底极板开关断开后，时变的采样噪声即转化为固定电压，行为特性与常规失调电压高度相似。文献 [2] 利用 SAR ADC 固有的环路反馈特性，构建了采样噪声消除机制，在最终输出数字码中实现噪声的有效剥离。

然而，文献 [2] 和 [3] 均指出该技术存在关键失效因素——\textbf{饱和问题}。在高增益条件下，噪声采样时间窗口受到严格限制；而在低增益条件下，对后级电容噪声及比较器噪声的抑制效果显著减弱。这一\textbf{增益--时间窗口权衡}问题成为制约技术性能的核心瓶颈。

针对上述局限，后续研究提出了多种改进方案：

\begin{table}[htbp]\centering\caption{改进方案对比}\label{tab:improve}\small
\begin{tabular}{lll}\toprule
\textbf{改进方案} & \textbf{特点} & \textbf{局限性} \\\midrule
带宽耦合分析 [4] & 与跨导参数耦合缓解饱和 & 仍存在设计约束 \\
反馈放大器噪声消除 [5] & 改善能效 & 缺失预放大级对后级噪声与失调的抑制功能 \\
预采样电容技术 [7], [8] & 缓解饱和问题 & 电路复杂度显著上升；Split-sampling 存在电荷泄露 \\\bottomrule
\end{tabular}
\end{table}

深入分析可知，上述技术瓶颈的根源在于\textbf{输出失调存储（OOS）技术的固有局限性}：失调存储电容置于信号链路中，限制了高增益放大器的应用，且易导致 SNR 下降及输出饱和。

\textbf{电流域技术的突破}

为突破传统 OOS 技术的限制，UC Irvine 团队提出了基于电流域失调消除的新型架构 [6]。该技术在 180 nm CMOS 工艺下验证了其优越性：
\begin{itemize}
  \item 失调电压存储电容未置于信号链路，高增益放大器使用不受限
  \item 失调探测阶段呈现低增益特性，对饱和具有更强鲁棒性
  \item 建立时间更短，饱和问题通常不构成主要设计约束
\end{itemize}

本研究基于上述技术背景，系统阐述电流域失调消除技术的核心原理，完成与传统电压域架构的对比实验验证，并在 PVT 及瞬态噪声仿真环境下进行全面的性能评估。\textbf{Phase II 阶段进一步推进至 10-MS/s 全面表征，采用 NF=10 噪声注入双 $t_d$ 方法首次完成了 CAAZ 消除效率 $\eta$ 的全 PVT 定量表征，并通过双批次独立 CSV 导出实现了跨批次验证。}

\subsection{传统电压域技术的局限性}

\begin{prinbox}{原理解析}
OOS 架构的核心矛盾：预放大器增益 $A$ 扮演双重角色——高增益抑制后级噪声，但导致输出摆幅增大引发饱和；低增益避免饱和，但对次级噪声衰减不足。文献被迫采用 $A \approx 6$ 的折中设计。
\end{prinbox}

传统电压域噪声消除技术采用级联输出失调存储（OOS）架构，其核心思想是将失调和噪声信息存储在输出端电容上，随后在转换阶段予以抵消。然而，该架构存在根本性的\textbf{增益--饱和死循环}，严重制约其性能上限。

\textbf{增益-饱和矛盾的理论分析}

\begin{table}[htbp]\centering\caption{增益条件对比}\label{tab:gain}\small
\begin{tabular}{llll}\toprule
\textbf{增益条件} & \textbf{优势} & \textbf{劣势} & \textbf{设计约束} \\\midrule
高增益 ($A > 10$) & 后级噪声抑制能力强 & 输出摆幅易超电源轨 & 饱和风险显著增加 \\
低增益 ($A < 6$) & 饱和裕度充足 & 次级噪声衰减不足 & 难以满足高精度需求 \\\bottomrule
\end{tabular}
\end{table}

\textbf{典型妥协方案}：文献 [2] 被迫采用 $A \approx 6$ 的折中设计，试图在饱和裕度和噪声抑制之间寻求平衡。然而，该增益值对次级噪声和失调的衰减程度有限（约 15.6 dB），难以满足高精度应用需求。

\textbf{残余噪声公式推导}

传统 OOS 架构的等效输入残余噪声可表示为：
\begin{equation}
v_{ns,res} = v_{ns1} e^{-2\Delta t/\tau} + \frac{v_{ns,AZ}}{A} + \frac{v_{ns,comp}}{A}
\label{eq:oos}
\end{equation}

\begin{table}[htbp]\centering\caption{残余噪声分量释义}\label{tab:noisecomp}\small
\begin{tabular}{lll}\toprule
\textbf{噪声项} & \textbf{表达式} & \textbf{物理意义} \\\midrule
主采样噪声残余 & $v_{ns1} e^{-2\Delta t/\tau}$ & 衰减依赖于建立时间常数 $\tau$ 和时间窗口 $\Delta t$ \\
自归零电容噪声 & $v_{ns,AZ} / A$ & 衰减因子为 $1/A$ \\
比较器噪声 & $v_{ns,comp} / A$ & 受增益 $A$ 抑制 \\\bottomrule
\end{tabular}
\end{table}

该公式表明，传统架构需要在建立时间和增益之间进行精细权衡，设计空间极为受限。

\textbf{时间窗口约束}

除增益约束外，时间窗口 $\Delta T_{az}$ 的选择同样面临两难境地。时间窗口过短会导致噪声混叠问题加剧，削弱噪声消除效果；时间窗口过长则增加饱和风险，同时降低整体转换速度。文献 [3] 详细分析了时间窗口与噪声混叠的关系，指出当 $\Delta T_{az} < 3\tau$ 时，噪声混叠将显著恶化信噪比。

\subsection{研究目标}

\begin{findbox}{核心亮点}
\textbf{三大核心目标}：① 在 180nm 工艺下实现 12-bit SAR ADC；② 系统验证电流域架构的抗饱和与噪声抑制优势；③ 实现系统级功耗优化并完成跨批次验证。
\end{findbox}

\textbf{Phase I（3--4月）：10-MS/s 方案验证}

\begin{table}[htbp]\centering\caption{Phase I 目标}\label{tab:phase1-goal}\small
\begin{tabular}{lll}\toprule
\textbf{指标} & \textbf{目标值} & \textbf{说明} \\\midrule
分辨率 & 12-bit & 理想信噪比 74 dB \\
采样率 & 10 MS/s & 奈奎斯特带宽 5 MHz \\
工艺 & 180 nm CMOS & 成熟工艺，成本效益高 \\\bottomrule
\end{tabular}
\end{table}

\textbf{Phase II（4--5月）：10-MS/s 全面表征}

\begin{table}[htbp]\centering\caption{Phase II 目标与达成}\label{tab:phase2-goal}\small
\begin{tabular}{llll}\toprule
\textbf{指标} & \textbf{目标值} & \textbf{实测 (TT)} & \textbf{结论} \\\midrule
ENOB (NF=1) & $>$11 bit & \textbf{11.61 bit} & $\checkmark$ 超额 \\
SNR (NF=1) & $>$70 dB & \textbf{71.81 dB} & $\checkmark$ \\
SFDR & $>$75 dB & \textbf{84.01 dB} & $\checkmark$ 超额 \\
CAAZ $\eta$ & $>$85\% & \textbf{87.9\%} & $\checkmark$ \\
总功耗 & $<$1 mW & \textbf{720 $\mu$W} & $\checkmark$ \\
PVT 覆盖 & 5 角 & TT/FF/SS/SF/FS & $\checkmark$ \\\bottomrule
\end{tabular}
\end{table}

\textbf{电流域架构优势验证}
\begin{itemize}
  \item \textbf{抗饱和验证}：Phase I 通过对比实验验证 CAAZ 在高增益条件下的稳定性优势，可用增益 $A > 50$
  \item \textbf{噪声抑制验证}：Phase I 完成噪声分解（KT/C 压缩 5.53$\times$）；Phase II 通过 NF=10 双 $t_d$ 方法定量表征 $\eta$
  \item \textbf{PVT 鲁棒性验证}：两阶段均在 TT、FF、SS、SF、FS 全工艺角及全温度范围内验证
\end{itemize}

\textbf{系统级功耗优化}

Phase II 实现了全芯片功耗 720 $\mu$W（模拟 65\%/数字 35\%），预放大器 PREWOCM 占 41\%，为当前最大优化空间。通过噪声消除技术，显著降低了对采样电容的要求，从而减轻外部驱动负担，实现整体系统能效的提升。

\subsection{创新点预览}

\begin{findbox}{核心亮点}
\textbf{四大创新点}：① 电容负担转移——从高阻抗外部节点转移到低阻抗内部节点；② 电流域自动调零——打破增益--饱和耦合；③ 成熟工艺高能效实现——降低制造成本；④ \textbf{NF=10 $\eta$ 定量表征}——首次通过噪声注入方法系统量化 CAAZ 消除效率。
\end{findbox}

\textbf{电容负担转移机制}：将大电容从高阻抗外部节点转移到低阻抗内部节点。CAAZ 架构将噪声存储电容 $C_{AZ}$（300 fF）置于内部低阻抗节点（$1/g_{mn}$），建立时间常数显著降低两个数量级，从根本上缓解了饱和问题。

\textbf{电流域自动调零技术}：
\begin{itemize}
  \item 噪声采样阶段采用低增益配置（$A_{samp} < 10$）以避免饱和
  \item SAR 转换阶段释放高增益（$A > 50$）以抑制后级噪声
  \item 实现了增益与饱和的解耦
\end{itemize}

\textbf{NF=10 $\eta$ 定量表征（Phase II 新增）}：
\begin{itemize}
  \item 100$\times$噪声注入（NF=10），双 $t_d$ 差分方法
  \item 首次完成五角 PVT CAAZ $\eta$ 全面覆盖：$\eta$=79.7\%--91.2\%
  \item 双批次跨验证确保结论稳健（偏差 $<$1.4 p.p.）
\end{itemize}

\textbf{成熟工艺高能效实现}：在 180 nm 成熟工艺下实现 10-MS/s 设计。全芯片功耗仅 720 $\mu$W。

\section{CAAZ 技术原理深度解析}
\label{sec:principle}

\begin{findbox}{核心亮点}
CAAZ 架构通过\textbf{三相工作时序}实现噪声提取与转换：采样阶段低增益避免饱和，转换阶段释放高增益抑制后级噪声。核心创新在于将噪声存储电容置于低阻抗节点，实现增益与饱和的解耦。
\end{findbox}

\subsection{CAAZ 电流域架构原理}

CAAZ 架构通过三相工作时序，将噪声提取转移到电流域，从根本上解决了增益--饱和矛盾 [6]。

\textbf{三相工作时序}

\begin{table}[htbp]\centering\caption{三相工作时序}\label{tab:phases}\small
\begin{tabular}{llll}\toprule
\textbf{阶段} & \textbf{工作状态} & \textbf{增益} & \textbf{说明} \\\midrule
主采样 & $V_{in}$ 与 $kT/C$ 噪声被冻结在 $C_1$ & — & 采样开关断开 \\
电流域噪声采样 & 负载管二极管连接 & $A_{samp} < 10$ & 低增益避免饱和 \\
SAR 转换 & 负载管切换为电流源 & $A > 50$ & 高增益抑制噪声 \\\bottomrule
\end{tabular}
\end{table}

\begin{figure}[htbp]\centering
\includegraphics[width=0.55\textwidth]{Fig6_Preamplifier.png}
\caption{电流域噪声消除架构}\label{fig:preamp}
\end{figure}

\textbf{电流域噪声消除原理推导}

\textbf{关键节点电压}

\begin{table}[htbp]\centering\caption{CDS 操作关键节点电压}\label{tab:cds}\small
\begin{tabular}{lp{7.5cm}}\toprule
\textbf{相位} & \textbf{节点电压表达式} \\\midrule
采样相位 ($t_1$) & $V_{C1} = V_{in}(t_1) + v_{ns1}$ \\
放大相位 ($t_2$) & $V_{C2} = g_{mp}R_{out}[V_{in}(t_2) - V_{in}(t_1) - v_{ns1} - v_{os}] + v_{ns2}$ \\\bottomrule
\end{tabular}
\end{table}

\textbf{电流域噪声消除机制}

1. \textbf{差分电流产生}：输入管的跨导 $g_{mp}$ 将电压差转换为电流
   $$I_{diff} = g_{mp} [V_{in}(t_2) - V_{in}(t_1) - v_{ns1} - v_{os}]$$

2. \textbf{噪声电流注入}：负载管跨导 $g_{mn}$ 将 $t_2$ 时刻噪声转换为电流
   $$I_{ns} = g_{mn} v_{ns2}$$

3. \textbf{总电流变化}：
   $$\Delta I = g_{mp}[g_{mn}R_{out} - 1][V_{in}(t_2) - V_{in}(t_1) - v_{ns1} - v_{os}] + g_{mn}v_{ns2}$$

\begin{impbox}{噪声消除原理}
该表达式清晰地表明，输出电压与 $t_1$ 时刻采样的噪声 $v_{ns1}$ 相关。在理想情况下，如果 $v_{ns1}$ 和 $v_{ns2}$ 是相关的，通过巧妙的设计可使它们在输出端被抵消，从而实现噪声消除。
\end{impbox}

\textbf{等效输入噪声转换}

通过输出阻抗 $R_{out}$ 将电流转换回电压：
$$\Delta V = R_{out} \Delta I = g_{mp}R_{out}[g_{mn}R_{out} - 1][V_{in}(t_2) - V_{in}(t_1) - v_{ns1} - v_{os}] + g_{mn}R_{out}v_{ns2}$$

等效输入噪声可简化为：
$$v_{ns,eq} \approx \frac{v_{ns1}}{g_{mp}R_{OA}} + \frac{v_{ns2}}{g_{mp}R_{OA}} \cdot \frac{g_{mn}}{g_{mp}}$$

该式表明，主噪声 $v_{ns1}$ 和次级噪声 $v_{ns2}$ 均被增益 $g_{mp}R_{OA}$ 衰减。

\begin{prinbox}{关键公式推导}
\textbf{残余主噪声}：$v_{ns1,res} = v_{ns1}/(g_{mp}R_{OA})$，衰减由晶体管本征参数决定，不依赖建立时间。\textbf{次级噪声}：$v_{ns,AZ\_in} \approx \sqrt{kT/C_{AZ}}$，无法被增益抑制，需通过增大 $C_{AZ}$ 降低。
\end{prinbox}

\textbf{残余主噪声}：
$$v_{ns1,res} = \frac{v_{ns1}}{g_{mp}R_{OA}}$$

与传统架构的 $e^{-2\Delta t/\tau}$ 相比，该衰减由晶体管本征参数决定，不依赖建立时间。180-nm CMOS 中 $g_{mp}R_{OA} > 50$，衰减超过 34 dB。

\textbf{次级噪声}：
$$v_{ns,AZ\_in} = \sqrt{\frac{kT}{C_{AZ}}} \cdot \frac{g_{mn}}{g_{mp}} \approx \sqrt{\frac{kT}{C_{AZ}}}$$

由于 Phase 2 时 $g_{mn} \approx g_{mp}$，次级噪声无法被增益抑制，必须通过增大 $C_{AZ}$ 直接降低。\textbf{Phase II 优化结果}：$C_{AZ}=300$ fF 已满足需求，500 fF 无额外收益。

\subsection{电容负担转移机制}

\begin{prinbox}{核心创新}
传统架构电容置于高阻抗输出端（$\tau \sim 100\,\mathrm{k\Omega} \cdot C$），CAAZ 将 $C_{AZ}$ 置于低阻抗节点（$\tau \sim 1/g_{mn} \cdot C_{AZ}$），时间常数降低两个数量级。
\end{prinbox}

\begin{table}[htbp]\centering\caption{电容位置对比}\label{tab:cap}\small
\begin{tabular}{llll}\toprule
\textbf{架构} & \textbf{电容位置} & \textbf{节点阻抗} & \textbf{特性} \\\midrule
传统 OOS & 预放大器输出端 & $\sim$100 k$\Omega$ & 易导致输出饱和 \\
CAAZ & 内部电流节点 & $\sim$1 k$\Omega$ & 增益与饱和解耦 \\\bottomrule
\end{tabular}
\end{table}

驱动时间常数对比：
$$\tau_{external} = R_{driver} C_1 \gg \tau_{internal} = \frac{C_{AZ}}{g_{mn}}$$

\textbf{Phase II 验证}：$C_{AZ}=300$ fF 在 $t_d=26$ ns 窗口内充分建立（$\eta_{TT}=87.9\%$）；$t_d=24.05$ ns 时 $\eta\approx0$。

\section{实验平台搭建}
\label{sec:platform}

\begin{findbox}{核心亮点}
基于 TSMC 180nm CMOS 工艺，核心模块包括：CDAC（$C_{unit}=5$ fF）、CAAZ 预放大器（$g_{mp}R_{OA}>50$）、自举开关、StrongARM 锁存器。Phase II 推进至 10-MS/s 全面表征，扩展 PVT 至五角全覆盖。
\end{findbox}

\subsection{工艺与设计环境}

\begin{table}[htbp]\centering\caption{设计环境}\label{tab:env}\small
\begin{tabular}{lll}\toprule
\textbf{项目} & \textbf{配置} & \textbf{说明} \\\midrule
工艺节点 & 180 nm CMOS & 成熟工艺，成本效益高 \\
电源电压 & 1.8 V & 标准数字电源电压 \\
EDA 工具 & Cadence Virtuoso + Spectre 23.1 (APS) & 业界标准模拟设计流程 \\
仿真类型 & Transient + Noise, 15 $\mu$s & 瞬态噪声仿真 \\
分析工具 & MATLAB R2024b / Python ADCToolbox & FFT 后处理与数据分析 \\\bottomrule
\end{tabular}
\end{table}

\subsection{核心电路模块}

\begin{table}[htbp]\centering\caption{核心电路模块演进}\label{tab:blocks}\small
\begin{tabular}{lp{2.2cm}p{2.2cm}p{3cm}}\toprule
\textbf{模块} & \textbf{Phase I} & \textbf{Phase II} & \textbf{设计要点} \\\midrule
CDAC & $C_{unit}=5$ fF & 同 Phase I & Split-DAC 结构，冗余位设计 \\
CAAZ 预放大器 & $g_{mp}R_{OA}>20$ & $g_{mp}R_{OA}>50$ & pMOS 输入对，nMOS 有源负载 \\
CAAZ 电容 & $C_{AZ}=500$ fF & $\mathbf{300}$ \textbf{fF} & Phase II 优化：300 fF 已足够 \\
$t_d$ 优化 & 24.1--25 ns & \textbf{26 ns（最优）} & Phase II 通过扫参确定 \\
自举开关 & $R_{on}<400\,\Omega$ & 同 Phase I & 电荷注入补偿，底极板采样 \\
StrongARM 锁存器 & 再生时间 $<$0.3 ns & 同 Phase I & Kickback 噪声隔离 \\
异步时钟生成 & 26.6--35 ns 可调 & 26 ns 固定 & PRST/RST 相位控制 \\\bottomrule
\end{tabular}
\end{table}

\subsection{仿真测试平台}

\textbf{Phase I：瞬态噪声仿真配置（10-MS/s）}

\begin{table}[htbp]\centering\caption{Phase I 仿真配置}\label{tab:sim1}\small
\begin{tabular}{lll}\toprule
\textbf{参数} & \textbf{设置值} & \textbf{说明} \\\midrule
仿真时长 & 15 $\mu$s & 覆盖足够多采样周期 \\
瞬态噪声带宽 & $f_s/2 = 5$ MHz & 奈奎斯特带宽 \\
工艺角 & TT/FF/SS/SF/FS & 全工艺角覆盖 \\
温度范围 & $-40^\circ$C 至 $125^\circ$C & 典型工作范围 \\\bottomrule
\end{tabular}
\end{table}

\textbf{Phase II：CAAZ 效率表征（10-MS/s）}

\begin{table}[htbp]\centering\caption{Phase II 仿真配置}\label{tab:sim2}\small
\begin{tabular}{lll}\toprule
\textbf{参数} & \textbf{设置值} & \textbf{说明} \\\midrule
仿真时长 & 15 $\mu$s & 瞬态噪声 \\
信号频率 & 4.765625 MHz & 61/128 相干采样 \\
输入幅度 & 890 mV$_{pp-diff}$ & $-0.5$ dBFS \\
NF 配置 & 1（基线）/ 10（100$\times$注入） & CAAZ 效率表征 \\
$t_d$ 扫描 & 24.05 / 26.0 ns & 双 $t_d$ $\eta$ 计算 \\
工艺角 & TT/FF/SS/SF/FS & 五角全覆盖 \\
噪声注入开关 & I1（时钟链），I14（采样开关） & — \\
批次验证 & 旧导出 / 新导出 & 跨批次交叉验证 \\\bottomrule
\end{tabular}
\end{table}

\section{核心实验发现与验证}
\label{sec:results}

\begin{findbox}{核心亮点}
Phase I（10-MS/s）：KT/C 噪声功率压缩 5.53$\times$，噪声占比从 31.40\% 降至 7.60\%。Phase II（10-MS/s 全面表征）：ENOB=11.61 bit（TT），$\eta$=87.9\%，全五角 PVT 覆盖，双批次交叉验证。
\end{findbox}

\subsection{Phase I：噪声分解与 KT/C 抑制量化（10-MS/s）}

实验聚焦于瞬态噪声开启下的频谱分析。通过 FFT 后处理提取各工艺角动态性能。下述数据基于 $C_{AZ}=500$ fF，$t_d=25$ ns。

\textbf{各工艺角性能汇总}

\begin{table}[htbp]\centering\caption{PVT 瞬态噪声仿真结果 ($t_d=25$ ns)}\label{tab:pvt1}\small
\begin{tabular}{lcccc}\toprule
\textbf{工艺角} & \textbf{温度} & \textbf{SNDR (dB)} & \textbf{ENOB (bit)} & \textbf{备注} \\\midrule
TT & 27$^\circ$C & 69.29 & 10.94 & 典型情况 \\
FF & 25$^\circ$C & 68.90 & 10.76 & 快工艺角 \\
SS & 0$^\circ$C & 68.28 & 11.01 & 慢工艺角 \\
SF$_{-40}$ & $-40^\circ$C & 69.36 & 10.99 & 慢N快P（低温） \\
SF$_{35}$ & 35$^\circ$C & 68.12 & 10.97 & 慢N快P（常温） \\
FS & $-13^\circ$C & 68.62 & 10.98 & 快NMOS慢PMOS \\\bottomrule
\end{tabular}
\end{table}

\textbf{噪声抑制效果对比（$t_d = 24.1$ ns）}

\begin{table}[htbp]\centering\caption{有/无 CAAZ 抑制效果对比}\label{tab:supp}\small
\begin{tabular}{lccc}\toprule
\textbf{工艺角} & \textbf{有抑制 ENOB (bit)} & \textbf{无抑制 ENOB (bit)} & \textbf{抑制增益 (bit)} \\\midrule
TT & 10.78 & 10.35 & $+$0.43 \\
FF & 10.64 & 10.36 & $+$0.28 \\
SS & 10.68 & 10.50 & $+$0.18 \\
SF$_{-40}$ & 10.92 & 10.51 & $+$0.41 \\
SF$_{35}$ & 10.59 & 10.12 & $+$0.47 \\
FS & 10.77 & 10.37 & $+$0.40 \\\bottomrule
\end{tabular}
\end{table}

\begin{figure}[htbp]\centering
\includegraphics[width=0.48\textwidth]{Fig4_Noise_Suppression.pdf}\hfill
\includegraphics[width=0.48\textwidth]{Fig5_Noise_Waterfall.pdf}
\caption{噪声抑制效果对比（左）与噪声源剥离瀑布图（右）}\label{fig:noise}
\end{figure}

\subsection{Phase I：噪声抑制效果定量分析}

基于 12-bit SAR ADC 的瞬态噪声仿真，对 KT/C 采样噪声抑制效果进行 100\% 复核无误的绝对功率域定量分析。

\textbf{固定已知条件（完全按实测数据，无任何篡改）}

\begin{table}[htbp]\centering\caption{四工况实验条件}\label{tab:conditions}\small
\begin{tabular}{clcl}\toprule
\textbf{工况} & \textbf{工况描述} & \textbf{实测 SNR} & \textbf{核心说明} \\\midrule
1 & 仅紧量化噪声 & 74.0 dB & 仅量化噪声，无其他任何噪声源 \\
2 & 量化 + 原始无抑制 KTC 噪声 & 71.4 dB & 仅量化$+$原始 KTC 噪声 \\
3 & 量化 + KTC Cancellation 抑制技术 & 73.4 dB & 仅量化$+$抑制后残余 KTC 噪声 \\
4 & 整机全链路实测 & 70.8 dB & 完整系统开 KTC 抑制的真实结果 \\\bottomrule
\end{tabular}
\end{table}

\textbf{核心公式}（归一化信号功率 $S=1$，行业通用标准）：
$$\text{SNR(dB)} = 10\lg(S/P_{noise}) \implies P_{noise} = 10^{-\text{SNR(dB)}/10}$$

非相关噪声总功率 = 各独立噪声功率\textbf{线性直接叠加}（非 dB 叠加）。

\textbf{步骤 1：基础噪声功率精确计算}

\begin{table}[htbp]\centering\caption{噪声功率基础值}\label{tab:pnoise}\small
\begin{tabular}{lll}\toprule
\textbf{噪声项} & \textbf{精确功率值} & \textbf{计算依据} \\\midrule
纯量化噪声 $P_Q$ & $\boldsymbol{3.9811\times 10^{-8}}$ & $10^{-74.0/10}$ \\
量化$+$原始 KTC 总噪声 $P_{Q+K_{raw}}$ & $\boldsymbol{7.2444\times 10^{-8}}$ & $10^{-71.4/10}$ \\
量化$+$抑制后 KTC 总噪声 $P_{Q+K_{sup}}$ & $\boldsymbol{4.5709\times 10^{-8}}$ & $10^{-73.4/10}$ \\
整机开抑制总噪声 $P_{ALL_{sup}}$ & $\boldsymbol{8.3176\times 10^{-8}}$ & $10^{-70.8/10}$ \\\bottomrule
\end{tabular}
\end{table}

\textbf{步骤 2：独立噪声分量拆分（100\% 自洽无歧义）}

\begin{table}[htbp]\centering\caption{独立噪声分量拆分}\label{tab:split}\small
\begin{tabular}{lll}\toprule
\textbf{独立噪声分量} & \textbf{精确功率值} & \textbf{计算逻辑} \\\midrule
原始未抑制 KTC 噪声 $P_{K_{raw}}$ & $\boldsymbol{3.2633\times 10^{-8}}$ & $P_{Q+K_{raw}} - P_Q$ \\
抑制后残余 KTC 噪声 $P_{K_{sup}}$ & $\boldsymbol{5.8981\times 10^{-9}}$ & $P_{Q+K_{sup}} - P_Q$ \\
整机固定其他噪声 $P_{OTHER}$ & $\boldsymbol{3.7467\times 10^{-8}}$ & $P_{ALL_{sup}} - P_{Q+K_{sup}}$ \\\bottomrule
\end{tabular}
\end{table}

\textbf{步骤 3：核心技术指标（无任何估算，纯实测推导）}

\textbf{KTC 噪声本身的抑制能力}

\begin{table}[htbp]\centering\caption{KTC 抑制核心指标}\label{tab:ktc-core}\small
\begin{tabular}{lc}\toprule
\textbf{指标} & \textbf{数值} \\\midrule
KTC 噪声功率压缩倍数 & \textbf{5.53 倍} \\
KTC 噪声功率衰减 & \textbf{$-$7.43 dB} \\
KTC 噪声功率削减比例 & \textbf{81.93\%} \\
纯量化$+$KTC 场景下 SNR 提升 & \textbf{$+$2.0 dB} \\\bottomrule
\end{tabular}
\end{table}

\textbf{整机全链路真实性能提升（严格功率叠加计算）}

\begin{table}[htbp]\centering\caption{整机全链路 SNR 提升}\label{tab:full-sys}\small
\begin{tabular}{lc}\toprule
\textbf{指标} & \textbf{数值} \\\midrule
整机关 KTC 抑制总噪声 $P_{ALL_{raw}}$ & $P_Q + P_{K_{raw}} + P_{OTHER} = \boldsymbol{1.0991\times 10^{-7}}$ \\
整机关 KTC 抑制对应 SNR & $10\lg(1/P_{ALL_{raw}}) \approx \boldsymbol{69.59\ \mathrm{dB}}$ \\
整机开 KTC 抑制实测 SNR & $\boldsymbol{70.80\ \mathrm{dB}}$ \\
$\checkmark$ 整机全链路真实 SNR 提升 & \textbf{$+$1.21 dB} \\\bottomrule
\end{tabular}
\end{table}

\begin{figure}[htbp]\centering
\includegraphics[width=0.65\textwidth]{Fig8_Noise_Pie_Charts.pdf}
\caption{KTC Cancellation 技术对 ADC 整机噪声分布的优化效果}\label{fig:pie8}
\end{figure}

\begin{findbox}{Phase I 核心结论}
本次提出的 KTC Cancellation 采样噪声抑制技术，可将原始 KTC 采样噪声功率压缩 5.53 倍，实现 7.43 dB 的噪声衰减，KTC 噪声能量削减 81.93\%。在仅含量化与 KTC 噪声的理想模块级场景中，该技术可实现 2.0 dB 的 SNR 提升；在全链路整机场景中，实测系统 SNR 由 69.59 dB 提升至 70.80 dB（$+$1.21 dB）。噪声分布对比结果显示，该技术可将 KTC 噪声在整机总噪声中的占比从 29.69\% 大幅降低至 7.09\%，有效缓解了高精度 SAR ADC 中采样噪声对系统动态范围的限制，具备显著的工程应用价值。
\end{findbox}

\subsection{Phase II：10-MS/s NF=10 CAAZ 效率定量表征}

Phase II 推进至 10-MS/s 全面表征，采用 NF=10 噪声注入双 $t_d$ 方法首次完成 CAAZ 高效率 $\eta$ 的全 PVT 定量表征。

\textbf{NF=1 基线性能（$t_d=26$ ns）}

\begin{table}[htbp]\centering\caption{NF=1 基线全角性能}\label{tab:nf1-data}\small
\begin{tabular}{lcccc}\toprule
\textbf{Corner} & \textbf{温度} & \textbf{SNR (dB)} & \textbf{ENOB (bit)} & \textbf{SFDR (dB)} \\\midrule
TT & 27$^\circ$C & \textbf{71.81} & \textbf{11.61} & \textbf{84.01} \\
FF & $-40^\circ$C & 71.16 & 11.49 & 81.85 \\
SS & 85$^\circ$C & 71.48 & 11.55 & 83.86 \\
SF & 85$^\circ$C & 71.29 & 11.53 & 83.52 \\
FS & $-40^\circ$C & 70.92 & 11.47 & 81.50 \\\bottomrule
\end{tabular}
\end{table}

\textbf{关键发现}：全五角 PVT ENOB 均 $>$11.4 bit，12-bit ADC 裕量充足。ENOB 极差仅 \textbf{0.14 bit}（11.47--11.61），CAAZ 全角稳定。SFDR 全部 $>$81 dB，线性度良好。

\textbf{$t_d$ 优化：24.05 ns vs 26.0 ns}

\begin{table}[htbp]\centering\caption{$t_d$ 延展对 SNR 的影响}\label{tab:td-comp}\small
\begin{tabular}{lccc}\toprule
\textbf{Corner} & \textbf{SNR$_{24}$ (dB)} & \textbf{SNR$_{26}$ (dB)} & $\Delta$SNR (dB) \\\midrule
TT & 70.39 & 71.81 & \textbf{$+$1.42} \\
FF & 70.79 & 71.16 & $+$0.37 \\
SS & 69.66 & 71.48 & \textbf{$+$1.82} \\
SF & 69.44 & 71.29 & \textbf{$+$1.85} \\
FS & 70.30 & 70.92 & $+$0.62 \\\bottomrule
\end{tabular}
\end{table}

\begin{figure}[htbp]\centering
\includegraphics[width=0.85\textwidth]{Fig1_PVT_Performance.pdf}
\caption{PVT 全角性能：ENOB（左）与 SNR（右），$t_d=24.05$ vs 26.0 ns}\label{fig:pvt2}
\end{figure}

\textbf{结论}：高温角（SS/SF, 85$^\circ$C）受益最大（SNR $+$1.8--1.9 dB），低温角（FF/FS, $-40^\circ$C）几乎无变化（$g_m$ 大, 24ns 已充分建立）。\textbf{26.0 ns 为最优设计点}，26.5 ns 无额外收益。

\textbf{CAAZ 消除效率 $\eta$}

通过 NF=10（100$\times$噪声注入）双 $t_d$ 方法计算。具体采用新导出 CSV 数据：
$$\eta = 1 - \frac{101 \cdot 10^{-\Delta\text{SNR}/10} - 1}{100},\quad \Delta\text{SNR} = \text{SNR}_{26\text{ns}} - \text{SNR}_{24.05\text{ns}}$$

\begin{table}[htbp]\centering\caption{CAAZ $\eta$ 全角表征（新导出数据）}\label{tab:eta-full}\small
\begin{tabular}{lccccc}\toprule
\textbf{Corner} & \textbf{SNR$_{26}$ (dB)} & \textbf{SNR$_{24}$ (dB)} & $\eta$ (新导出) & $\eta$ (旧导出) & \textbf{降噪倍} \\\midrule
TT & 64.10 & 55.22 & \textbf{87.9\%} & 86.8\% & \textbf{8.3$\times$} \\
FF & 66.40 & 57.01 & 89.4\% & 89.4\% & 9.4$\times$ \\
SS & 61.03 & 53.33 & 83.9\% & 82.5\% & 6.2$\times$ \\
SF & 61.60 & 54.84 & \textbf{79.7\%} & 78.8\% & 4.9$\times$ \\
FS & 67.06 & 56.94 & \textbf{91.2\%} & 90.9\% & \textbf{11.4$\times$} \\\bottomrule
\end{tabular}
\end{table}

\begin{figure}[htbp]\centering
\includegraphics[width=0.65\textwidth]{Fig2_CAAZ_Eta.pdf}
\caption{CAAZ $\eta$ 跨批次对比（两批次偏差 $<$1.4 p.p.）}\label{fig:eta2}
\end{figure}

\textbf{关键发现}：
\begin{itemize}
  \item TT 角 $\eta=\textbf{87.9\%}$，满足 $>$85\% 目标
  \item FS 角 $\eta=\textbf{91.2\%}$ 全 PVT 最优（低温 $+$ fast-N, $g_{mn}$ 最大）
  \item SF 角 $\eta=\textbf{79.7\%}$ 全 PVT 最差（85$^\circ$C $+$ slow-N fast-P, $g_{mn}$ 最不足）
  \item 双批次 $\eta$ 偏差 $<$1.4 p.p. $\rightarrow$ \textbf{CAAZ 效率结论跨批次稳健}
\end{itemize}

\textbf{工艺角依赖性物理机制}：CAAZ 建立时间常数 $\tau_{CAAZ}\approx C_{AZ}/g_{mn}$。FF/FS 角（$-40^\circ$C）$g_{mn}$ 最大，建立最充分（$\eta>89\%$）。SF/SS 角（85$^\circ$C）$g_{mn}$ 降低，建立不足。SF 角 slow-N fast-P 组合使 NMOS 负载 $g_{mn}$ 最弱，$\eta$ 最低。

\textbf{功耗分解 (TT, NF=1, $t_d=26$ ns)}

\begin{table}[htbp]\centering\caption{功耗分解}\label{tab:power2}\small
\begin{tabular}{lcc}\toprule
\textbf{模块} & \textbf{功耗 ($\mu$W)} & \textbf{占比} \\\midrule
COMPOWER（预放大器$+$SR$+$驱动器） & 363.7 & 50\% \\
~~PREWOCM & 298.4 & 41\% \\
~~SRPOWER & 30.93 & 4\% \\
SWITCHPOWER（开关网络） & 97.14 & 14\% \\
LOGICPOWER（SAR\_LOGIC） & 149.0 & 21\% \\
SYNCPOWER（SYNC） & 110.4 & 15\% \\\midrule
\textbf{总计} & \textbf{720.3} & 100\% \\\bottomrule
\end{tabular}
\end{table}

模拟部分占 $\sim$65\%，数字部分占 $\sim$35\%。\textbf{预放大器 PREWOCM 是最大功耗源（41\%），为当前最大优化空间。}

\begin{figure}[htbp]\centering
\includegraphics[width=0.75\textwidth]{Fig3_Power_Breakdown.pdf}
\caption{功耗分解图}\label{fig:power}
\end{figure}

\textbf{跨批次数据验证}

\begin{table}[htbp]\centering\caption{双批次跨验证核心指标}\label{tab:cross2}\small
\begin{tabular}{lccc}\toprule
\textbf{指标} & \textbf{旧导出} & \textbf{新导出} & \textbf{偏差} \\\midrule
COMPOWER (TT) & 363.2 $\mu$W & 365.2 $\mu$W & $+$2.0 $\mu$W \\
$\eta$ (TT) & 86.8\% & 87.9\% & $+$1.1 p.p. \\
ENOB (TT, NF=10, 26ns) & 10.31 & 10.31 & 0 \\
SNR (TT, NF=10, 26ns) & 63.68 dB & 64.10 dB & $+$0.42 dB \\
SWITCHPOWER & 97.2 $\mu$W & 110.7 $\mu$W & \textbf{$+$14\%} \\\bottomrule
\end{tabular}
\end{table}

COMPOWER、$\eta$、ENOB、SNR 等核心指标跨批次高度一致。SWITCHPOWER 存在系统性 $+$14\% 偏差，经核查系第二次 CSV 导出时 Calculator 表达式不同所致，非电路实际变化。\textbf{CAAZ 效率定性结论跨批次稳健。}

\subsection{失调注入测试}

\begin{prinbox}{原理解析}
仅开启采样噪声时 ENOB 达 11.59 bit，证明 CAAZ 架构配合极小采样电容可近乎完美消除 $kT/C$ 噪声。引入后级全噪声后 ENOB 回落至 11.08 bit，说明比较器噪声是当前精度瓶颈。
\end{prinbox}

\textbf{噪声源与失调电压对 ADC 性能的影响分解 ($t_d = 25$ ns)}

\begin{table}[htbp]\centering\caption{噪声/失调分解分析}\label{tab:offset}\small
\begin{tabular}{lccr}\toprule
\textbf{仿真条件配置} & \textbf{ENOB 范围 (bit)} & \textbf{SNR 范围 (dB)} & \textbf{平均 ENOB} \\\midrule
仅开启采样噪声 & 11.42--11.70 & 70.42--72.54 & 11.59 bit \\
全噪声开启 (0 mV 失调) & 10.94--11.18 & 67.51--69.28 & 11.08 bit \\
全噪声开启 (2 mV 失调) & 11.03--11.25 & 68.12--69.36 & 11.13 bit \\
全噪声开启 (5 mV 失调) & 10.85--10.92 & 67.37--67.76 & 10.89 bit \\\bottomrule
\end{tabular}
\end{table}

分析表明，在仅开启前端采样噪声时，各工艺角下的平均 ENOB 高达 11.59 bit，峰值达到 11.70 bit。这证明了本文提出的 CAAZ 架构配合 5 fF 极小采样电容，能够近乎完美地消除 $kT/C$ 噪声的限制。当引入比较器等后级全系统热噪声时，ENOB 回落至 11.08 bit，说明比较器噪声是当前制约系统精度的主要瓶颈。

\textbf{失调注入测试结果（隔离寄生电容，$t_d = 25$ ns）}

\begin{table}[htbp]\centering\caption{失调注入测试}\label{tab:offset2}\small
\begin{tabular}{lccc}\toprule
\textbf{注入失调 (mV)} & \textbf{ENOB 范围 (bit)} & \textbf{平均 ENOB (bit)} & \textbf{退化幅度} \\\midrule
0 & 10.94--11.18 & 11.08 & 基准 \\
2 & 11.03--11.25 & 11.13 & $\sim$0 \\
5 & 10.85--10.92 & 10.89 & $-0.19$ \\\bottomrule
\end{tabular}
\end{table}

\begin{figure}[htbp]\centering
\includegraphics[width=0.48\textwidth]{Fig6_Offset_Sensitivity.pdf}\hfill
\includegraphics[width=0.48\textwidth]{Fig7_Time_Window_Sweep.pdf}
\caption{失调敏感度分析（左）与时间窗口扫参（右）}\label{fig:offset-sweep}
\end{figure}

\begin{findbox}{失调鲁棒性结论}
注入失调从 0 mV 增至 5 mV 时，ENOB 仅下降约 0.19 bit。CAAZ 架构对等效输入失调展现出强鲁棒性，验证了电流域消除机制对失调的有效处理能力。
\end{findbox}

\subsection{Phase I $\rightarrow$ Phase II 演进分析}

\begin{table}[htbp]\centering\caption{两阶段关键指标对比}\label{tab:evolution2}\small
\begin{tabular}{lcc}\toprule
\textbf{维度} & \textbf{Phase I (10-MS/s)} & \textbf{Phase II (10-MS/s)} \\\midrule
ENOB (TT, NF=1) & 10.94 bit & \textbf{11.61 bit} ($+$0.67) \\
SNR (TT) & 67.51 dB & \textbf{71.81 dB} ($+$4.30) \\
$C_{AZ}$ & 500 fF & \textbf{300 fF} ($-$40\%) \\
$t_d$ & 25 ns & \textbf{26 ns}（扫参最优） \\
噪声表征方法 & 四工况分解 & \textbf{NF=10 $\eta$ 法}（更系统化） \\
PVT 覆盖 & 6 角（含 SS\_mim） & \textbf{5 角（含 SF）} \\
跨批次验证 & 无 & \textbf{双批次}（$\eta$ 偏差 $<$1.4 p.p.） \\
功耗数据 & 比较器 272 $\mu$W & \textbf{全芯片 720 $\mu$W} \\
SFDR & N/A & \textbf{84.01 dB} \\\bottomrule
\end{tabular}
\end{table}

Phase II 通过优化 $C_{AZ}$、建立 NF=10 $\eta$ 表征方法论，实现了从验证架构可行到全面定量表征的跨越。双批次跨验证为 CAAZ 技术的数据可靠性提供了有力支撑。

\section{技术优势与结论}
\label{sec:conclusion}

\subsection{技术优势总结}

\begin{findbox}{三大核心优势}
\textbf{抗饱和}：增益与饱和解耦，可用增益 $A > 50$（$8\times$ 传统 OOS）。\textbf{噪声抑制}：主噪声衰减因子 $\approx$0.02，$\eta_{TT}=87.9\%$，不依赖建立时间。\textbf{系统能效}：小采样电容降低前端驱动负担，全芯片 720 $\mu$W。
\end{findbox}

\textbf{抗饱和优势}

\begin{table}[htbp]\centering\caption{抗饱和对比}\label{tab:anti-sat}\small
\begin{tabular}{lcc}\toprule
\textbf{参数} & \textbf{传统 OOS 架构} & \textbf{CAAZ 架构} \\\midrule
可用增益范围 & $A < 10$ & $A > 50$ \\
增益-饱和关系 & 强耦合 & 解耦 \\
时间窗口灵活性 & 受饱和限制 & 仅受建立时间限制 \\\bottomrule
\end{tabular}
\end{table}

\textbf{噪声抑制优势}

\begin{table}[htbp]\centering\caption{噪声抑制机制对比}\label{tab:noise-adv}\small
\begin{tabular}{lccc}\toprule
\textbf{架构} & \textbf{衰减机制} & \textbf{衰减因子} & \textbf{典型值} \\\midrule
传统 OOS & 建立时间衰减 & $e^{-2\Delta t/\tau}$ & $\sim$0.1 \\
CAAZ & 增益衰减 & $1/(g_{mp}R_{OA})$ & $\approx$0.02 \\\bottomrule
\end{tabular}
\end{table}

\subsection{后续研究计划}

\begin{warnbox}{待解决挑战}
\textbf{关键技术挑战}：① 预放功耗优化——PREWOCM 占 41\%，目标缩小至 200 $\mu$W 以下；② SF 角改善——$\eta=79.7\%$ 为最差角，需增大 $t_d$ 余量或提升预放 $g_{mn}$；③ 版图后仿——PEX 反标验证；④ 低频噪声——NF=1 时 N=5 退化 $\sim$0.3 bit。
\end{warnbox}

\begin{table}[htbp]\centering\caption{优先行动清单}\label{tab:actions}\small\setlength{\tabcolsep}{6pt}
\begin{tabular}{c>{\raggedright}p{3cm}>{\raggedright}p{3.5cm}>{\raggedright}p{5.5cm}}\toprule
\textbf{级} & \textbf{项目} & \textbf{目标} & \textbf{说明} \\\midrule
P0 & SWITCHPOWER 表达式 & 统一 Cadence 定义 & 跨批次 $+$14\% 根因确认 \\
P1 & 预放电流优化 & 总功耗 $<$600 $\mu$W & PREWOCM 占 41\%, 最大优化空间 \\
P1 & SF 角改善 & $\eta_{SF}>85\%$ & 增大 $t_d$ 或提升预放 $g_{mn}$ \\
P1 & CCLK H-tree & skew $<$20 ps & 版图对称布线验证 \\
P2 & PEX 寄生反标 & 后仿验证 & 版图完成后必须执行 \\
P2 & 低频 1/f 噪声排查 & 闪烁噪声确认 & N=5 退化约 0.3 bit \\\bottomrule
\end{tabular}
\end{table}

\subsection{总结}

本报告系统研究了基于电流调整自动调零（CAAZ）技术的 SAR ADC 设计。

\textbf{Phase I}（3--4月）：在 10-MS/s 下完成 CAAZ 架构原理验证，通过四工况瞬态噪声仿真进行噪声分解。KTC 噪声功率压缩 5.53$\times$（81.93\% 能量削减），系统 SNR 提升 1.21 dB。各工艺角 ENOB 均达 10.7--11.0 bit。

\textbf{Phase II}（4--5月）：推进至 10-MS/s 全面表征，首次通过 NF=10 双 $t_d$ 方法完成全 PVT $\eta$ 定量表征（79.7\%--91.2\%），双批次跨验证确认结论稳健。基线性能达 ENOB=11.61 bit, SNR=71.81 dB, SFDR=84.01 dB（TT 角），全五角 ENOB $>$11.4 bit，总功耗 720 $\mu$W。SF 角（85$^\circ$C, slow-N fast-P）为当前设计瓶颈（$\eta=79.7\%$），后续可通过预放电流优化和 $t_d$ 余量改善进一步提升。

\begin{findbox}{核心创新点汇总}
\begin{tabular}{lll}\toprule
\textbf{创新点} & \textbf{技术贡献} & \textbf{量化意义} \\\midrule
电容负担转移 & 大电容从高阻抗$\rightarrow$低阻抗节点 & $\tau$ 降低两个数量级 \\
电流域自动调零 & 电流域噪声提取，打破增益-饱和 & $A > 50$（传统 $A < 6$） \\
NF=10 $\eta$ 表征 & 100$\times$ 噪声注入双 $t_d$ 定量方法 & $\eta_{TT}=87.9\%$ 首次全 PVT \\
双批次跨验证 & 独立 CSV 导出交叉校验 & $\eta$ 偏差 $<$1.4 p.p. \\
成熟工艺高能效 & 180nm 工艺下 10-MS/s & 720 $\mu$W, ENOB 11.61 bit \\\bottomrule
\end{tabular}
\end{findbox}

\vspace{1cm}
\begin{center}\rule{0.5\textwidth}{0.4pt}
{\small Rev.4.0 (2026-05-13): 忠实保留原始框架，仅更新数据值}

{\small Rev.1.0 (2026-04-18): 初版}
\rule{0.5\textwidth}{0.4pt}
\end{center}

\end{document}
